% 本文件由 LaTeX 排版 Agent 根据有效 translation.json 自动生成。
% author=Jerome; work_id=3010; title=驳维吉兰提乌斯

\chapter{驳维吉兰提乌斯}

% structure: p0001 source=h1 target=omitted reason=page-title-already-rendered-by-parent

关于维吉兰提乌斯的详细情况，即本论著（一夜之功的成果）所抨击的对象，可查阅1844年由达勒姆议事司铎\Person{吉利}博士出版的\WorkTitle{维吉兰提乌斯及其时代}一书。然而，或许简要说明一下，有助于读者理解：他大约生于370年，在\Place{卡拉克里斯}（靠近\Place{康韦纳}，即\Place{科曼日}），这是从\Place{阿基坦}通往\Place{西班牙}的罗马道路上的一个驿站。他的父亲可能是客栈老板，\Person{维吉兰提乌斯}似乎自幼便从事父亲的行业。他生性好学，教父史家\Person{苏尔皮基乌斯·塞维鲁斯}在那片地区拥有田产，便收留了他，并可能让他管理自己的庄园。他被祝圣后，经由\Person{诺拉的保利努斯}（他是\Person{苏尔皮基乌斯·塞维鲁斯}的朋友）介绍，认识了\Person{耶柔米}（当时住在\Place{伯利恒}，时为395年）。与\Person{耶柔米}共处相当长一段时间后，他请求离开，并匆忙离去，未作任何解释。回到\Place{高卢}后，他定居于家乡。\Person{耶柔米}听说他散布谣言，说他倾向于\Person{奥力振}的观点，并以其他方式诽谤他和他的朋友，便写了一封严厉的斥责信（\WorkTitle{书信第六十一封}）。\Person{维吉兰提乌斯}的那部促使\Person{耶柔米}写下以下论著的著作，写于公元406年；并非\Quote{仓促之间，出于他离开\Place{伯利恒}时可能感到的愤怒}，而是过了六七年之后。他所抨击的迷信行为包括：（1）对圣人遗骸的敬礼，即用贵重容器或丝巾包裹着遗骸在教堂内巡行供人亲吻，并向亡者祈祷；（2）深夜在殉道者大殿守夜，随之而来的丑闻，点燃无数蜡烛，所谓奇迹等；（3）向\Place{耶路撒冷}捐赠施舍，\Person{维吉兰提乌斯}主张这些施舍最好花在每个教区的穷人身上，以及修道的贫穷誓愿；（4）对贞洁的过分推崇。\par

该教区的主教，\Person{图卢兹的埃克苏佩里乌斯}，强烈支持\Person{维吉兰提乌斯}的观点，这些观点开始广泛传播。司铎\Person{里帕里乌斯}向\Person{耶柔米}提出了控诉，\Person{耶柔米}立即表示愤慨，并提出如果\Person{维吉兰提乌斯}的著作被送来，他将详细作答。406年，他通过\Person{西西尼乌斯}收到了这部著作，\Person{西西尼乌斯}当时正带着施舍前往\Place{东方}。有确切的评论说，这篇论著比\Person{耶柔米}所写的任何其他作品都更缺乏理性，而更多是谩骂。尽管如此，作者当时受到了主要教会人士的追随，\Person{维吉兰提乌斯}所抨击的做法几乎毫无阻碍地盛行，直至十六世纪。\par

1. 世界产生了许多怪物；在\BibleBook{依撒意亚}{先知书}中，我们读到半人马和塞壬、仓鸮和鹈鹕。\Person{约伯}以奥秘的语言描述了里外雅堂和贝希摩斯；刻耳柏洛斯和\Place{斯廷法罗斯湖}的怪鸟、\Place{厄律曼托斯}的野猪和\Place{涅墨亚}的狮子、喀迈拉和九头蛇许德拉，都是诗话中的传说。\Person{维吉尔}描述了卡库斯。\Place{西班牙}产生了三身革律翁。只有\Place{高卢}没有怪物，却一直出产勇敢和雄辩之人。突然之间，\Person{维吉兰提乌斯}——更准确地说，是\OriginalTerm{瞌睡虫}{Dormitantius}——出现了，被不洁之灵驱使，对抗基督之神，否认殉道者坟墓应受宗教敬礼。他说，守夜应当谴责；除了复活节，决不可唱\Quote{阿肋路亚}；节制是异端；洁德是色欲的温床。据说，\Person{欧福耳布斯}在\Person{毕达哥拉斯}身上重生，同样，\Person{约维尼亚努斯}的腐败思想在这个人身上重现；因此，在他身上，我们与在他的前辈身上一样，必然遇见魔鬼的陷阱。这句话可恰当地用在它身上：\Quote{邪恶的种子，因你们父亲的罪恶，准备你们的儿女被屠杀。}\Person{约维尼亚努斯}被\Place{罗马}教会的权威所谴责，在雉鸡和猪肉之间，吐出了——不如说打嗝出了——他的灵魂。而现在，这位\Place{卡拉克里斯}的客栈老板，就本村的名字而言，是一位\Person{昆体良}，只是沉默而不善言辞，却往酒里掺水。他利用自己熟知的花招，企图把他背信弃义的毒液掺入公教信仰中；他攻击童贞，憎恨洁德；他与世俗之徒狂欢，却谴责圣人的斋戒；他杯觥交错卖弄哲学，以甜美的圣咏旋律来抚慰自己，同时咂嘴品尝干酪饼；除非在宴席上，他不屑于聆听\Person{大卫}、\Person{耶杜顿}、\Person{阿撒夫}和\Person{科辣黑}子孙的诗歌。我诉说出这些，与其说是取乐，不如说是悲伤，因为我无法自制，对宗徒和殉道者所受的冤屈充耳不闻。\par

2. 说来可耻，据说有些主教竟与他同流合污——若他们尚可称为主教——他们只祝圣那些先结了婚的人为执事；他们不相信独身者有贞洁——不但如此，他们更藉着对众人怀有恶意的猜疑，清楚地表明自己所能声称的圣洁生活程度如何；除非晋秩候选人带着怀孕的妻子、怀中抱着啼哭的婴儿来到他们面前，他们就不授予基督的圣事。教会在东方将如何行事？埃及的教会和那些属于宗徒之座的教会又将如何？它们只接纳童男或守贞者，或已婚却放弃夫妇权利的人为圣职人员。这就是\Person{Dormitantius}的教导：他放纵情欲，鼓励肉欲，使青年时期本已沸腾的肉欲加倍炽热，或更确切地说，通过与妇女的交合来熄灭它；以致我们与猪毫无区别，与禽兽或马也无分别，经上记着说：\BibleRef{耶 5:8} \BibleQuote{他们像发情的马匹向女人狂奔；各人都向他邻人的妻子嘶叫。} 这就是圣神藉达味的口所说的：\Quote{不要像马和骡那样无知。} 又论及\Person{Dormitantius}和他的朋友们：\Quote{用辔头和嚼环勒住那不愿接近你的下巴。}\par

3. 但现在我们该引述他自己的话，并详细回答他。因为，也许出于恶意，他会再次曲解我，说我编造了一个案例，以展示我与之争论的修辞和演讲才能，就像我写给高卢的那封关于一对母女不和的信一样。我现在口述的这篇短文，是献给可敬的长老\Person{Riparius}和\Person{Desiderius}的，他们来信说他们的堂区因靠近他而受玷污，并托我们的弟兄\Person{Sisinnius}寄来他醉醺醺时吐出的那些书卷。他们还说，有些人因倾向他的恶习而赞同他的亵渎之词。他无论在言语还是知识上都是个野蛮人。他的文体粗糙。他甚至不能捍卫真理；但为了平信徒和那些罪孽深重、常学不学、总不达真理之知的可怜妇女们，我愿意花一夜的功夫来处理他那些可悲的琐碎言论，否则我将显得藐视那些恳请我从事这项工作的可敬人士的来信。\par

4. 他确实很好地代表了他的种族。他出身于一群强盗和从各处聚集起来的人（我指的是\Person{Cn. Pompey}征服西班牙后，在匆忙返回举行凯旋式时，从比利牛斯山带下来并聚集到一个城镇的那些人，因此该城得名\Place{Convenæ}），他通过攻击天主的教会，延续了他们强盗的行径。如同他的祖先\Person{Vectones}、\Person{Arrabaci}和\Person{Celtiberians}一样，他袭击高卢的教会，不是高举十字架的旗帜，而是魔鬼的旗帜。\Person{Pompey}在东方也做了同样的事。在征服基利家与伊索里亚的海盗和强盗之后，他在基利家与伊索里亚之间建立了一座以自己名字命名的城市。然而，那座城市至今遵守其祖先的规矩，没有出现任何\Person{Dormitantius}；但高卢却供养着一个本土的敌人，并看着一个失心疯、应当被绑上\Person{Hippocrates}所推荐的束缚带的人坐在教会中。在他的其他亵渎之词中，可以听到他说：\Quote{你们不仅尊敬，且几乎崇拜你们用小器皿带在身边的那个东西（无论它是什么），那有什么必要？} 又在同一本书中说：\Quote{你们为什么亲吻和崇拜一块用布包着的粉末？} 又在同一本书中说：\Quote{在宗教的外衣下，我们看到一种几乎像是异教徒的仪式被引入教会：太阳还在照耀时，就点燃成堆的蜡烛，而到处都有人在亲吻和崇拜一块用贵重布料包裹的廉价粉末。这些人真是对真福殉道者致以极大的敬意，竟以为他们需要靠劣质的蜡烛来增添光荣，而事实上，那在中座中的羔羊，以祂全部威仪的光辉，已为他们照亮了。}\par

5. 疯子！世上谁曾敬拜过殉道者？谁曾以为人是天主？难道\BibleRef{宗 14:11}保禄和巴尔纳伯，当吕考尼雅人以为他们是尤比特和墨丘利，想要向他们献祭时，不是撕裂衣服，声明他们不过是人吗？这并非因为他们不比尤比特和墨丘利更优秀——后者不过是早已死去的人——而是因为在外邦人的错误观念下，原应归于天主的尊荣正被给予他们。关于伯多禄，我们也读到同样的事：当科尔乃略想要敬拜他时，他用手扶起他来，并说\BibleRef{宗 10:26}\Quote{起来，因为我也只是一个人。}你竟敢提到\Quote{你们带在一个小容器里并敬拜的那个神秘东西或别的什么？}我想知道你所称为\Quote{东西或别的什么}的是什么。请你更清楚地告诉我们（好让你的亵渎之言不受约束）你所说的\Quote{用贵重布料包裹在小容器里的一丁点粉末}是什么意思。这无非是殉道者的圣髑，他愤愤不平地看到它们被贵重的纱幔覆盖，而不是用破布或粗毛布包裹，或扔在粪堆上，好让维吉兰提独自身在醉醺醺的睡梦中被敬拜。难道我们进入宗徒大殿时，就是犯了亵圣之罪？君士坦提乌斯一世皇帝将安德肋、路加和弟茂德的圣髑移往君士坦丁堡时，难道犯了亵圣之罪？在这些圣髑面前，魔鬼喊叫，而住在维吉兰提内的魔鬼也承认他们感受到了圣人的影响。到了今天，阿卡狄乌斯皇帝在如此长时间后，将真福撒慕尔的骨骸从犹大运往色雷斯，难道犯了亵圣之罪？难道所有主教不仅要被视为亵圣者，还要被视为愚昧，因为他们携带那最无价值的东西——尘埃和骨灰——包在丝绸里、放在金器中？难道所有教会的人民都是傻瓜，因为他们前去迎接圣髑，以如同看见一位活生生的先知在当中的喜悦欢迎它们，以致从巴勒斯坦到加采东的一大群人异口同声地回响着对基督的赞美？他们其实是在敬拜撒慕尔，而不是基督——撒慕尔曾是基督的肋未人和先知。你表现不信，因为你只想到尸体，因而亵渎。请阅读福音\BibleRef{玛 22:32}\Quote{亚巴郎的天主，依撒格的天主，雅各伯的天主：他不是死人的天主，而是活人的天主。}既然他们是活着的，那么，用你的话说，他们就不是被锁在光荣的牢狱中。\par

6. 你说，宗徒和殉道者的灵魂，或居住在亚巴郎的怀抱中，或安息之所，或天主的祭台底下，且不能离开自己的坟墓，随意临在。他们似乎属于元老等级，不受最恶劣的监狱和杀人犯的陪伴，而是被隔离在福岛和乐土中，享受自由而体面的监禁。你要为天主立法吗？你要把宗徒们锁起来吗？以致他们被拘禁直到审判之日，不与他们的主同在，尽管论及他们记载着：\BibleRef{默 14:4}\BibleQuote{他们跟随羔羊，无论祂往哪里去。}如果羔羊无处不在，那么对与羔羊同在的人，也必须相信同样的事。魔鬼和恶魔遍游世界，且以极快的速度随处显现；殉道者在流了血之后，却被关在棺材里，无法逃脱，这合理吗？你在你的小册子中说，当我们活着时，可以彼此代祷；但一旦死了，无人能为他人祈祷而被俯听，尤其因为殉道者虽然\BibleRef{默 6:10}呼求为他们血的伸冤，却从未能获得所求。如果宗徒和殉道者还在肉身时能为他人祈祷——那时他们仍当为自己挂虑——那么一旦他们赢得了冠冕，得胜凯旋，岂不更该如此？一个普通人梅瑟，常为六十万武装之众从天主求得宽恕；\BibleRef{宗 7:59}斯德望，主的追随者和首位基督徒殉道者，为迫害他的人求赦免；而一旦他们开始与基督同在的生活，难道权力反不如从前？保禄宗徒\BibleRef{宗 27:37}说，有二百七十六个人在船中赐给了他；而当他离世开始与基督同在后，难道他必须闭口无言，不能为那些普世相信他福音的人说一句话吗？活的狗维吉兰提乌斯比死的狮子保禄更强吗？如果我承认保禄在灵性上死了，我就可以按训道篇那样说。事实上，圣人不是称为死了，而是被说成睡了。因此\BibleRef{若 11:11}拉匝禄，将要复活，被说是睡着了。宗徒\BibleRef{得前 4:13}禁止得撒洛尼人为那些睡了的人忧愁。至于你，醒着时是睡着的，写作时也是睡着的，你给我拿出一本伪经，以厄斯德拉之名被你和你的同类诵读，书中写着死后无人敢为他人祈祷。我从未读过那书：何必采用教会不接受的著作？你总不能拿巴耳撒慕、巴耳贝鲁、摩尼的宝库、以及可笑的名字留西波拉斯来反驳我；不过，可能因为你住在比利牛斯山脚下，毗邻伊比利亚，你跟随古代异端者巴西里德及其所谓知识的难以置信的奇谈，宣扬被普世权威所谴责的东西。我这样说，是因为在你的短论中，你引撒罗满好像他支持你，尽管撒罗满从未写过所述的那些话；所以，既然你有第二厄斯德拉，你也会有第二撒罗满。你若愿意，可以阅读所有圣祖和先知虚构的启示，学会后，在妇女的织布工坊中歌唱它们，或者更便当地在你们的酒馆中朗读，以这些悲伤的小调刺激无知的群众添满他们的酒杯。\par

7. 至于蜡烛的问题，我们并不像你徒劳地歪曲的那样在白天点蜡烛，而是用它们的慰藉来照亮黑夜的黑暗，等待黎明，免得我们像你一样瞎眼，在黑暗中沉睡。如果有某些无知单纯的在俗信友，或者热心的妇女——关于她们我们可以真确地说：\BibleRef{罗 10:2}\BibleQuote{我承认他们对天主有热忱，但不是按照真知}——为了尊敬殉道者而采纳这习惯，这对你造成什么损害？从前连宗徒们也曾抱怨香膏浪费了，但他们被主的声音责备。基督不需要香膏，殉道者也不需要蜡烛的光；但那女人倾注香膏以尊敬基督，她内心的虔诚被接受了。凡是点这些蜡烛的人，按照他们的信德获得赏报，正如宗徒所说：\Quote{各人当在自己的心意上充盈。}你称这等人为拜偶像者吗？我不否认，我们所有信基督的人，都已脱离拜偶像的错谬。因为我们不是生而为基督徒，而是藉着重生成为基督徒。难道因为我们从前敬拜偶像，现在就不该敬拜天主，以免我们对祂和对偶像似乎表示同样的尊敬吗？前者是对偶像的尊敬，所以该憎恶那仪式；后者是对殉道者的敬礼，所以该容许同样的仪式。在整个东方教会，即使没有殉道者的圣髑，每当要宣读福音时都点蜡烛，尽管晨曦可能已染红天空，当然不是为了驱散黑暗，而是作为表达我们喜乐的方式。\BibleRef{玛 25:1}因此福音中的童女总是点着灯。宗徒们也被\BibleRef{路 12:35}嘱咐要束上腰，手里点着灯。关于若翰洗者，我们读到：\BibleRef{若 5:35}\BibleQuote{他是那燃烧并照耀的灯。}这样，以有形之光的形象，所代表的那光正是我们在圣咏中读到的：\BibleQuote{主，你的言语是我脚步的明灯，是我路径上的光明。}\par

8. 难道罗马主教做错了，当他向主献祭，在可敬的亡者伯多禄和保禄的骨骸之上——按我们所说，但按你所说，是在一堆毫无价值的尘土之上——并认为他们的坟墓配得上作为基督的祭台？不单是一位城市的主教犯错，而是全世界的主教都错了，他们不顾酒保\Person{维吉兰修斯}，进入亡者的大殿，其中\Quote{一块毫无价值的尘土和灰烬用布包裹着}，玷污自身并玷污一切。因此，按你所说，这些神圣建筑就像法利塞人的坟墓，外面粉饰洁白，里面却有污秽的遗骸，充满恶臭和不洁。然后他竟敢在这事上喷吐污秽，说道：\Quote{难道殉道者的灵魂爱他们的骨灰，环绕着它们，总在那里，唯恐若有人来祈祷而它们不在，它们就听不见吗？}哦，怪物，本当被放逐到地极！你嘲笑殉道者的遗骸，并与这异端之父\Person{欧诺米}一起，诽谤基督的教会？你不害怕与之为伍，并说出他对教会所说的同样的话来反对我们吗？因为所有他的追随者都拒绝进入宗徒和殉道者的大殿，以便他们可以敬拜死去的\Person{欧诺米}——他们认为他的书比福音更有权威——并相信真理之光在他身上，正如其他异端者坚持认为护慰者降在\Person{蒙丹}身上，并说摩尼本人就是护慰者。你甚至不能以自己是新型邪恶的发明者而自夸，因为你的异端早已爆发反对教会。然而，它遇到了一位对手\Person{戴尔都良}，一位非常博学的人，他写了一本著名的论著，恰当地称为\WorkTitle{蝎子毒}，因为正如蝎子像弓一样弯曲以施放毒刺，以前称为加音异端的，现在将毒液注入教会身体；它沉睡了，更确切地说是被埋葬了很长时间，但现在被\Person{多尔米坦修斯}唤醒了。我很惊讶你没有告诉我们，绝不可有殉道，因为天主不要求山羊和公牛的血，更不要求人的血。这就是你所说的，或者，即使你没有说，你也被认为是在断言它。因为你坚持要践踏殉道者的遗骸，你就禁止流他们的血，认为不值得尊敬。\par

9. 关于守夜和在殉道者大殿中频繁举行夜课，我在大约两年前写给可敬的司铎\Person{里帕里}的另一封信中已作了简短答复。你主张应该弃绝它们，免得我们似乎经常庆祝复活节，并且没有遵守惯常的年度守夜。如果是这样，那么就不应在主日向基督献祭，免得我们频繁庆祝主复活的复活节，并引入许多复活节代替一个的习惯。然而，我们不可将青年和卑劣妇女的过错和错误归咎于虔诚的人，这些过错和错误往往在夜间被发现。诚然，即使在复活节守夜中，通常也会暴露类似的事情；但少数人的过错并不能构成反对整个宗教的论据，而这些人不守夜也可能在自己的家里或别人的家里犯错。犹达斯的背信弃义并没有取消宗徒的忠诚。如果别人守夜守得不好，我们的守夜也不应因此被停止；相反，让那些为了满足情欲而睡觉的人被迫守夜，以便他们能保持贞洁。因为如果一件事做一次是好的，那么做多次就不可能坏；而如果有某种要避免的过错，错误不在于做得多，而在于根本上做了。因此，我们不应在复活节期间守夜，以免奸夫满足他们长期压抑的欲望，或者妻子有机会犯罪而不必担心丈夫的锁匙。那些很少出现的机会正是最被人渴望的。\par

10. 我无法详述可敬的长老书信中所包含的所有主题；我只从\Person{维吉兰提乌斯}的著作中提出几点。他反对在殉道者大殿中发生的表征和奇迹，并说这些是为不信的人服务的，而不是为信徒，好像现在的问题在于它们为谁的利益而发生，而不是凭借什么力量。就算表征属于不信的人，因为他们不肯听从言语和道理，便借着表征而相信。就连我们的主也为不信的人行了奇迹，然而我们主的奇迹并不因此受到非难，因为那些人是没有信德的；相反，它们应更值得钦佩，因为它们如此强大，以致连最刚硬的心也降服了，迫使人相信。所以，我不要你告诉我，表征是为了不信的人；但请你回答我的问题——为什么可怜而无价值的尘土和灰烬竟与这种表征和奇迹的神奇力量相关联？我看到，我看到，最不幸的凡人，你为何如此悲伤，以及是什么引起你的恐惧。那迫使你写下这些的不洁之神，常常因这无价值的尘土而受折磨，是的，此刻正在受折磨，尽管在你的身上他隐藏了伤口，但在别人身上他做了承认。你几乎会追随异教徒和不敬神的\Person{波菲利}和\Person{尤诺米乌}，假装这些是魔鬼的把戏，他们并不真的呼喊，而是假装痛苦。让我给你我的建议：去殉道者的大殿吧，有一天你会被洁净；你会在那里发现许多与你自己情况相似的人，而且会被点燃，不是被使你反感的殉道者的蜡烛，而是被无形的火焰；然后你会承认你现在所否认的，并会自由地宣告你的名字——那个以\Person{维吉兰提乌斯}的身份说话的你，实际上是\Person{墨丘利}，因为他贪财；或\Person{诺克图尔努斯}，按照\Person{普劳图斯}的\Quote{\WorkTitle{安菲特律翁}}，当\Person{朱庇特}连续两夜与\Person{阿尔克墨涅}私通并生下强大的\Person{赫拉克勒斯}时，他正在睡觉；或者至少是酒神\Person{巴克斯}，以醉酒闻名，酒壶挂在肩上，面孔永远红润，嘴唇起泡，言语放纵。\par

11. 有一次，当这个省份在半夜突然发生地震，把我们所有人都从睡梦中惊醒时，你，这个最谨慎最聪明的人，开始不穿衣服就祈祷，并让我们想起了\Person{亚当}和\Person{厄娃}在乐园中的故事；他们确实，当他们的眼睛睁开时感到羞耻，因为他们看到自己赤身露体，用树叶遮盖了自己的羞耻；但是你，被剥夺了衬衫和信德的你，在突如其来的恐惧中，带着昨晚的酒气，向弟兄们展示了你的智慧，以一种太过明显的方式暴露了你的裸体。这些就是教会的敌人；这些就是与殉道者之血争战的领袖；这就是抨击宗徒的演说家的样本，或者更确切地说，这些就是对着基督的门徒狂吠的疯狗。\par

12. 我承认我的恐惧，因为也许它会被认为是出于迷信。当我生气时，或心中有恶念时，或夜间的某种幻影迷惑了我时，我不敢进入殉道者的大殿，我浑身颤抖，身体和灵魂都是。你也许会微笑，并嘲笑说这与软弱妇人的妄念相当。如果是这样，我并不羞于拥有像那些首先看见复活的主的人那样的信德；他们被派往宗徒；他们，在我们的主和救主之母身上，被托付给了圣宗徒。如果你愿意，和世俗之人一起宣泄你的羞耻吧，我要和妇女一起禁食；是的，与那些面容见证其贞洁、因长期斋戒而面颊苍白、彰显基督的贞洁的虔诚男子一起。\par

十三、有些事似乎也令你困扰。你担心，如果节欲、清醒和斋戒在\Place{高卢}人当中扎根，你的酒馆就无利可图，你将无法彻夜进行你那属于魔鬼的守夜和狂饮作乐。此外，我从那些信函中得知，你藐视保禄的权威，甚至伯多禄、若望和雅各伯的权威——他们曾与保禄和巴尔纳伯伸出共融的右手，并吩咐他们记念穷人——你禁止向\Place{耶路撒冷}寄送任何金钱救济以惠及圣徒。现在，如果我对此回复，你立刻就会开口喊叫，说我是在为自己辩护。你果然对整个团体如此慷慨，以致若你没有来到耶路撒冷，并且挥霍你自己的钱或你赞助人的钱，我们都快要饿死了。我所说的是蒙福的宗徒保禄在他几乎所有的书信中所说的；他定下规矩，外邦人的教会应在每周的第一天，即主日，由每人捐款，然后送到耶路撒冷救济圣徒；要么由他自己的门徒送去，要么由他们亲自认可的人送去；如果认为合适，他本人会差人送去或亲自带去所收集的款项。同样在\WorkTitle{宗徒大事录}中，他对总督\Person{斐理斯}讲话时说：见：\BibleRef{宗 24:17} \Quote{过了多年，我上到耶路撒冷，是为了给我的民族带施舍和祭献，并履行我的誓愿；其间，他们在圣殿中看见我已行了洁净礼。}难道他不可以把从别人那里得来的礼物分配到世界其他地方，以及他用自己的信德所培育的初生教会吗？但他渴望施予圣地的穷人；他们为了基督而抛弃了自己的微小财产，全心转向事奉主。现在，如果我打算复述他书信中所有提倡和竭力主张把钱送到耶路撒冷和圣地给信徒的段落，那将花费太长时间；这不是为了满足贪婪，而是为了给予救济；不是为了积累财富，而是为了支撑贫弱身体的软弱，并抵御寒冷和饥饿。这习俗在\Place{犹太地}一直延续至今，不仅在我们中间，也在希伯来人中间，以致那些日夜默想上主的法律、在地上没有父亲、只有主为父的人，能藉着会堂和全世界的援助得到供养；为的是要达成均平\BibleRef{格后 8:14}——不是有些人得畅快，而其他人却困苦，而是以一些人的富余补足其他人的需要。\par

十四、你会回答说，每个人都可以在自己的家乡这样做，而且永远不乏需要靠教会资源供养的穷人。我们不否认，如果财力允许，救济品应分给所有穷人，甚至犹太人及\Place{撒玛黎雅}人。但宗徒教导说，施舍要给予众人\BibleRef{迦 6:10}，然而尤其是对信德之家的人。关于这些人，救主在福音中说：见：\BibleRef{路 16:9} \Quote{你们要藉着那不义的钱财结交朋友，使他们在你们进入永久的帐幕时，能接待你们。}什么！那些衣衫褴褛、污秽不堪、被狂怒的私欲所统治的可怜虫，他们现在和将来都一无所有，难道能拥有永久的居所吗？无疑，被称为有福的并非单纯穷人，而是神贫的人；正如经上所载：\Quote{用心照顾穷人和困苦人的，这人是有福的；在祸患的日子，上主必拯救他。}但事实上，在供养普通穷人时，所需要的不是心思，而是金钱。对于圣洁的穷人，心思确实有福的操练，因为你给予时，对方会脸红地接受，领受后又感到忧伤，他虽然播种了属灵的事物，却不得不收割你们的属肉之物。至于他的论证，即那些保留自己财产、并逐渐将增值部分分给穷人的人，比那些变卖一切、一次全部施舍的人更明智——对此，不是我在回答，而是主亲自作答：见：\BibleRef{玛 19:21} \Quote{你若愿意成为成全的，去变卖你所有的一切，施舍给穷人，然后来跟随我。}他是在对那愿意成为成全的人说话，这人像宗徒们一样，离开了父亲、船和网。你所赞同的人处于第二或第三等级；然而我们欢迎他，只要明白第一等优先于第二等，第二等优先于第三等。\par

十五、我还要说，我们的隐修士不应被你这毒舌和野蛮之咬而动摇决心。你关于他们的论证如下：如果所有人都隐居独处，那么谁去经常进出教堂？谁留下来赢取那些忙于世俗事务的人？谁能力劝罪人行德行之事？同样，如果所有人都像你一样愚蠢，谁能智慧？再顺着你的论证，童贞就不值得我们的认可了。因为如果所有人都守童贞，我们就没有婚姻；人类就会灭绝；婴儿就不会在摇篮里哭闹；助产士就会失业而沦为乞丐；而\Person{铎米当提乌}，独自一人、冻得蜷缩，会在床上失眠。事实上，德行是稀有之物，不为大众所热切追求。但愿所有人都像那少数人，经上说：\Quote{被召的人多，被选的人少。}监狱就会空荡荡的。但是，隐修士的职责并非教导，而是哀叹；或为自己，或为世界而悲伤，并且战战兢兢地期盼我主的来临。他深知自己的软弱和所携带器皿的脆弱，害怕绊倒，免得碰到什么东西而跌倒破碎。因此，他回避女性的视线，特别是年轻妇女，并且如此严格自律，以至于连安全的事都感到恐惧。\par

16. 你会说：为何要前往旷野？理由很明显：使我不听见、看不见你；使我不受你的狂妄所干扰；使我不与你交战；以免妓女的眼睛引诱我；以免美貌引我陷入非法的拥抱。你会回答：\Quote{这不是战斗，而是逃跑。你要站在战线上，穿上铠甲，抵抗你的仇敌，这样，你战胜后便可戴上冠冕。} 我承认我的软弱。我不愿冒着得胜的期盼去战斗，恐怕有时我会失去胜利。若我逃跑，我便躲避了刀剑；若我站立，我要么得胜，要么跌倒。但我何必放弃确定的事而追随不确定的事呢？我要么用盾牌，要么用双脚逃避死亡。你战斗的人，要么被击败，要么得胜。我逃跑的人，并不因逃跑而得胜；但我逃跑是为了确保我不被击败。若有一条蛇在你旁边，你睡觉并不安全。或许它不会咬我，但也许过一段时间它会咬我。我们称呼那些不比姐妹和女儿年长的女性为母亲，而我们并不愧于用虔敬的名义遮掩我们的恶习。一个隐修士在女子的住处做什么？那避开见证者的秘密谈话和眼神是什么意思？神圣的爱没有不安的欲望。再者，我们关于情欲所说的，也应应用于贪婪，以及一切借独处而避免的恶习。所以我们避开拥挤的城市，以免我们被迫去做那被催促去做的事，这与其说是出于本性，不如说是出于选择。\par

17. 正如我所说，应可敬的各位长老之请，我将这些论述的听写工作奉献了一个夜晚的辛劳，因为我的兄弟\Person{西西尼乌斯}正急忙启程前往埃及，他在那里要给圣人们提供救济，而且他急于离去。然而，若非如此，题目本身便如此明显地亵渎，以至于更需要作者的义愤，而非大量证据。但如果\Person{多尔米坦提乌斯}醒来再次辱骂我，并认为可以用那张同样亵渎的嘴——他用那张嘴攻击宗徒和殉道者——来诋毁我，我将为他花费比这短暂的夜间写作更多的时间。我将为他，以及为他的同伴们——无论是门徒还是导师——整夜守夜，这些人认为除非一个人结婚且他的妻子怀有身孕，否则无人配得上基督的职务。\par

\SourceRecord{Translated by W.H. Fremantle, G. Lewis and W.G. Martley.；Nicene and Post-Nicene Fathers, Second Series；Vol. 6.；Edited by Philip Schaff and Henry Wace.；Buffalo, NY: Christian Literature Publishing Co.,；1893.；<http://www.newadvent.org/fathers/3010.htm>.}
